\section{Discussion}

\subsection{Implementation reliability before physical refinement}

The main improvement arose from repairing the implementation rather than adding more detailed physical methods. Under the common diagnostic scenario, the original annual sensible demand of 2408.74 kWh fell to 131.76 kWh after repair, compared with 157.33 kWh in EnergyPlus. The accompanying reduction in hourly RMSE from 483.24 to 15.59 W indicates that the change affected agreement throughout the year, as well as the annual total. These findings support checking the accounting of gains, boundaries and heat-transfer pathways before interpreting differences as limitations of the simplified model formulation.

The diagnostic evidence makes the corrections physically interpretable. Repeated internal-gain additions inflated the gain total, while ground heat flow was calculated for an apartment with no ground contact. Removing these inconsistencies ensures that the model represents the intended apartment more faithfully. However, the decrease in internal gains cannot be equated directly with the decrease in heating or cooling demand. Thermostat operation and the timing of heat gains influence how these quantities are related. Most repairs were also evaluated together, so the annual improvement cannot be divided reliably among individual corrections. The infiltration ablation provides a separate annual effect, whereas the remaining checks establish correct behaviour without quantifying each repair's independent contribution.

The repaired implementation therefore provides a more defensible basis for further development, but it does not achieve equivalence with EnergyPlus. Its remaining annual underprediction of 16.25\% is material and should be retained when describing the model's performance. Energy-balance closure and successful software tests establish numerical consistency; they do not establish that the assumed building inputs or predicted loads are correct for an occupied dwelling.

\subsection{Physical detail and the choice of performance metric}

The angular-transmittance and exterior-convection comparisons show why annual and hourly performance must be assessed separately. Angular transmittance reduced cooling, whereas dynamic convection increased it under the tested conditions. Their combined effect on annual cooling was consequently smaller than either effect alone. This partial cancellation helps explain why a limited change in the annual total does not mean that the added methods were inactive.

With both methods enabled, total-load RMSE decreased from 15.59 to 15.02 W, while the magnitude of annual bias increased from 16.25\% to 17.06\%. The hourly improvement was also uneven: cooling RMSE decreased, but heating RMSE increased. The evidence therefore supports a modest improvement in the aggregate hourly comparison, rather than a general improvement across all performance measures. Retaining both methods in the Australian scenarios allows their joint behaviour to be assessed consistently; it should not be presented as proof that this combination is optimal.

The winter concentration of the heating shortfall provides a useful direction for further diagnosis. June--August accounted for 69.6\% of the repaired model's annual heating difference. This pattern identifies when disagreement is greatest, but does not isolate its cause. Further controlled comparisons could examine the treatment of neighbouring temperatures, heat-transfer coefficients, thermostat response and construction properties. Changing these factors together to match an annual total would obscure which mechanism explains the difference.

\subsection{Interpretation of the Australian adaptations}

The Australian adaptations primarily change the physical inputs represented by the model. Increasing the provisional leakage parameter from q$_{50}$ = 4 to 11 m$^3$/(h m$^2$) increased annual demand in both engines, but left their relative disagreement almost unchanged. This distinguishes the effect of representing a leakier envelope from an improvement in the calculation method. The construction-year lookup provides a reproducible way to select an assumed input; without measured permeability, it cannot establish that the selected value describes this apartment or the wider Australian housing stock.

Weather-dependent infiltration offers a way to examine the timing of leakage alongside outdoor conditions. At q$_{50}$ = 11, the prescribed weather-dependent flow increased heating by 0.92 kWh relative to constant flow, despite an annual mean multiplier below one. This result is consistent with the importance of when airflow occurs relative to heating operation. However, this comparison changes both the annual mean and the hourly pattern of airflow, so it does not isolate a pure timing effect. A further comparison using a constant flow equal to the weather-dependent annual mean would separate these influences more clearly.

The common airflow prescription was designed to control the comparison between engines. It uses a fixed 20 \textdegree C reference for the stack driver and the same hourly multiplier in both models. The results therefore assess responses to matched prescribed airflow, rather than validating a fully coupled airflow model driven by each engine's simulated indoor temperature. They also do not isolate a wind-only contribution or demonstrate performance under extreme wind conditions. The relatively small leakage effect is specific to the limited outdoor-exposed envelope and the prescribed continuous ventilation in this apartment; it does not establish that leakage is unimportant in other dwelling types.

The construction comparison produced a larger demand change than the tested leakage change, but its interpretation requires similar care. The transition from diagnostic zero-capacity construction to illustrative material-derived properties changes both resistance and heat storage. The reduction cannot be attributed to thermal mass alone or interpreted as the energy saving from a particular retrofit. The existing ISO capacity-class method was retained, so these runs also provide no evidence of improved representation of material layer order. Zero predicted cooling in the construction cases is a result of the assumed scenario and model response, rather than evidence that comparable apartments would never require cooling.

With both Australian inputs applied, pyBuildingEnergy predicted 76.51 kWh against 94.01 kWh in EnergyPlus, an underprediction of 18.61\%. Because each physical scenario has its own matched reference, these comparisons reveal how agreement changes with the inputs. They do not show a progressive improvement in accuracy from one scenario to the next. The combined case demonstrates that the regional inputs can be represented within the workflow while leaving a substantial discrepancy to be investigated.

\subsection{Scope and priorities for further evaluation}

The study provides a controlled comparison for one apartment and one weather dataset. Neighbouring temperatures, construction assemblies, glazing and ventilation remain assumed inputs, and the early stages deliberately use a diagnostic construction. These choices support the investigation of implementation behaviour but limit direct interpretation of the absolute loads as representative Australian residential demand. The reported quantities are sensible heating and cooling needs, not metered electricity use; equipment efficiency and a validated transient moisture model are outside the present comparison.

EnergyPlus serves as a computational reference rather than an observation of the building. Agreement with it is useful for identifying differences between implementations, but empirical validation requires measured conditions and demand. The present study also does not quantify uncertainty in the assumed inputs, assess peak-load accuracy separately, or establish a computational-speed advantage. Claims about operational prediction, equipment sizing or reliable real-time deployment would therefore exceed the evidence presented here.

Further work should prioritise separate ablations for the remaining implementation repairs, matched-mean airflow comparisons, and construction experiments that vary resistance and capacity independently. Measurements of permeability, neighbouring conditions and indoor temperatures would help constrain the case inputs. Evaluation across additional dwelling forms and climates would then test whether the remaining bias is systematic or specific to this apartment. Such comparisons should retain separate heating and cooling metrics, annual bias and hourly error, so that improvements in one measure do not conceal deterioration in another.

\section{Conclusion}

This study evaluated pyBuildingEnergy through implementation repairs, optional physical-method changes, and Australian leakage and construction inputs, using twelve annual simulations and seven matched EnergyPlus references. The largest improvement occurred after the implementation repairs: annual total bias changed from +1430.98\% to -16.25\%, and hourly total-load RMSE decreased from 483.24 to 15.59 W. Diagnostic checks supported the removal of repeated internal gains and inappropriate ground heat flow, while the corrected runs satisfied the numerical energy-balance criterion.

Additional physical detail produced mixed outcomes. Combining angular transmittance and dynamic exterior convection reduced total-load RMSE to 15.02 W but increased the magnitude of annual bias to 17.06\%. The Australian input comparisons showed sensitivity to both leakage and construction properties without resolving the discrepancy between engines. The combined configuration predicted 76.51 kWh of sensible demand against 94.01 kWh in EnergyPlus, remaining 18.61\% lower.

The contribution is a corrected, reproducible implementation and a controlled framework for distinguishing software repairs, physical-method effects and regional input assumptions. The results do not establish general validation for Australian housing or equivalence with EnergyPlus. The revised model provides a basis for further comparative research, with measured-input validation and broader case testing needed before stronger claims about predictive accuracy can be made.
